\section{Model Analysis}\label{sec:analysis}
\subsection{Value of the Complete Evidence Pathway}
Table~\ref{tab:core} compares \Model{} and the matched context-only model in dedicated paired cohorts: ETTh1$(.10,16)$ and PAMAP2$(.15,8)$. All three training-seed pairs favor \Model{} in both final MSE and MAE. The final MSE reductions are \GainETTMSE{} and \GainPAMMSE{}, while MAE reductions are \GainETTMAE{} and \GainPAMMAE{}, respectively. The coarse outputs already exhibit gains, and the differences remain after independently training the same refiner. This establishes an incremental value of explicit patch evidence beyond the matched contextual alternative.

\begin{table*}[t]
 \centering\small\setlength{\tabcolsep}{4pt}
 \begin{tabular}{llrrr}
\toprule
\Lang{数据集}{Dataset} & \Lang{模型}{Model} & \Lang{粗MSE}{\shortstack{Coarse\\MSE}} & \Lang{精MSE}{\shortstack{Refined\\MSE}} & \Lang{精MAE}{\shortstack{Refined\\MAE}} \\
\midrule
ETTh1 & \Model{} & $0.384873\,\pm\,0.006236$ & $0.334986\,\pm\,0.006672$ & $0.372039\,\pm\,0.004057$ \\
ETTh1 & \Lang{上下文对照}{Context-only} & $0.408733\,\pm\,0.001863$ & $0.353320\,\pm\,0.009562$ & $0.385618\,\pm\,0.010399$ \\
PAMAP2 & \Model{} & $0.336419\,\pm\,0.000680$ & $0.324572\,\pm\,0.001591$ & $0.230532\,\pm\,0.000622$ \\
PAMAP2 & \Lang{上下文对照}{Context-only} & $0.351320\,\pm\,0.002814$ & $0.332533\,\pm\,0.001960$ & $0.234430\,\pm\,0.001031$ \\
\bottomrule
\end{tabular}

 \caption{Paired evidence-pathway analysis. Both datasets use three training seeds and twenty shared test masks per model. The PAMAP2 mechanism cohort is separate from its formal benchmark cohort in Table~\ref{tab:formal}. Coarse and refined outputs come from the same selected model pair.}
 \label{tab:core}
\end{table*}

\subsection{Source Organization and Pointwise Reading}
Table~\ref{tab:components} shows that removing any one source increases mean final MSE in the two analyzed datasets. Cross-variable evidence has the largest source-removal MSE effect on ETTh1, whereas temporal evidence has the largest effect on PAMAP2. Their different contributions are consistent with source usefulness depending on the underlying signals.

Source-wise random selection and unified retrieval produce higher mean MSE than the complete system. Retaining provenance-specific organization therefore provides benefits beyond merely supplying the same number of legal fragments. Patch-shared reading changes the accuracy tradeoff: it gives lower ETTh1 MSE, while pointwise reading gives lower ETTh1 MAE and PAMAP2 MSE. We assess reading granularity using both error criteria.

\begin{table*}[t]
 \centering\small\setlength{\tabcolsep}{3pt}
 \begin{tabular}{lrrrr}
\toprule
\Lang{变体}{Variant} & ETTh1 MSE & ETTh1 MAE & PAMAP2 MSE & PAMAP2 MAE \\
\midrule
\Model{} & $0.334986\,\pm\,0.006672$ & $0.372039\,\pm\,0.004057$ & $0.322915$ & $0.230183$ \\
\Lang{去时间证据}{w/o temporal evidence} & $0.343612\,\pm\,0.006384$ & $0.381684\,\pm\,0.004346$ & $0.333096$ & $0.235096$ \\
\Lang{去跨变量证据}{w/o variable evidence} & $0.349493\,\pm\,0.018407$ & $0.381209\,\pm\,0.010604$ & $0.325663$ & $0.231132$ \\
\Lang{去历史证据}{w/o historical evidence} & $0.339364\,\pm\,0.004246$ & $0.376169\,\pm\,0.002895$ & $0.324277$ & $0.231290$ \\
\Lang{来源内随机}{Source-random} & $0.336177\,\pm\,0.006092$ & $0.377530\,\pm\,0.004732$ & $0.325863$ & $0.231077$ \\
\Lang{统一检索}{Unified retrieval} & $0.339484\,\pm\,0.005147$ & $0.380305\,\pm\,0.002030$ & $0.331974$ & $0.232928$ \\
\Lang{无候选容量匹配}{Context-only (matched)} & $0.353320\,\pm\,0.009562$ & $0.385618\,\pm\,0.010399$ & $0.334383$ & $0.235483$ \\
\Lang{Patch共享读取}{Patch-shared read} & $0.333639\,\pm\,0.004381$ & $0.374103\,\pm\,0.002174$ & $0.323958$ & $0.229965$ \\
\bottomrule
\end{tabular}

 \caption{Component and retrieval controls, after refinement. ETTh1 uses three training seeds (mean$\pm$sample SD); PAMAP2 component rows consistently use seed 2023. Every model is evaluated on twenty shared masks. Three-seed PAMAP2 results for \Model{} and the matched context-only model are given separately in Table~\ref{tab:core}.}
 \label{tab:components}
\end{table*}

\subsection{Observation-Conditioned Refinement}
For the paired \Model{} cohorts, refinement reduces MSE from 0.384873 to 0.334986 on ETTh1 and from 0.336419 to 0.324572 on PAMAP2, corresponding to 12.96\% and 3.52\%. The network combines temporal correction with reference information anchored to visible observations. Coordinate-level diagnostics separately account for error reductions and increases, making the aggregate correction traceable to its contributing positions.

The common-refiner study in Table~\ref{tab:backbone} further shows its effect across backbones. On BeijingAirQuality$(.20,8)$, T1 changes from 0.225523 to 0.129252 MSE after refinement; \Model{} changes from 0.179570 to 0.123941. The common correction stage accounts for a substantial part of the raw-output gap and preserves an additional advantage for \Model's coarse estimate. This comparison attributes the backbone and refinement effects separately.

\begin{table}[t]
 \centering\small\setlength{\tabcolsep}{3pt}
 \begin{tabular}{lrrrr}
\toprule
\Lang{骨干}{Backbone} & \Lang{粗MSE}{\shortstack{Coarse\\MSE}} & \Lang{精MSE}{\shortstack{Refined\\MSE}} & \Lang{粗MAE}{\shortstack{Coarse\\MAE}} & \Lang{精MAE}{\shortstack{Refined\\MAE}} \\
\midrule
T1 & $0.225523$ & $0.129252$ & $0.264340$ & $0.193203$ \\
\Lang{上下文对照}{Context-only} & $0.188083$ & $0.124760$ & $0.243393$ & $0.189678$ \\
\Model{} & $0.179570$ & $0.123941$ & $0.235067$ & $0.188432$ \\
\bottomrule
\end{tabular}

 \caption{A common refinement architecture applied to three frozen backbones on BeijingAirQuality$(.20,8)$, seed 2023, twenty test masks. Coarse and refined errors are evaluated on the same positions.}
 \label{tab:backbone}
\end{table}

\subsection{Missing-Condition Diagnostics}
We group evaluation coordinates by mask-defined gap length, nearest visible-point distance, patch visibility and window boundary. These groups are fixed before reading model outputs. On ETTh1$(.10,16)$, gains remain positive across all three training seeds for coordinates 17--32 and 33--95 steps from visible observations, supporting the use of patch evidence where immediate observations are sparse. For coordinates whose overlapping patches contain no visible observations, the external-to-patch context and other evidence sources still support recovery. The full stratified values and coordinate support are maintained with the numerical evidence tables.

\subsection{Accuracy and Cost}
Table~\ref{tab:efficiency} quantifies the additional cost of the candidate pathway. \Model{} and the matched context-only model have nearly equal effective parameter counts, while candidate construction and reading increase latency and allocated memory. On PAMAP2, the matched models use 575074 and 574794 effective parameters, with full-pipeline p50 latency 98.48 and 55.53ms per batch. This makes the measured accuracy gain assessable alongside the execution cost.

\begin{table}[t]
 \centering\small\setlength{\tabcolsep}{3pt}
 \begin{tabular}{llrrr}
\toprule
\Lang{数据集}{Dataset} & \Lang{模型}{Model} & \Lang{参数}{Params} & p50 (ms) & \Lang{显存}{VRAM} (GiB) \\
\midrule
PAMAP2 & \Model{} & 575,074 & 98.48 & 6.121 \\
PAMAP2 & \Lang{上下文对照}{Context-only} & 574,794 & 55.53 & 4.122 \\
ETTh1 & \Model{} & 353,248 & 17.63 & 1.067 \\
ETTh1 & \Lang{上下文对照}{Context-only} & 352,968 & 11.51 & 0.709 \\
\bottomrule
\end{tabular}

 \caption{Inference cost of selected seed 2023 models on the same exclusive RTX5090, FP32, batch 128. Median latency follows twenty warm-ups and one hundred synchronized measurements. VRAM denotes peak allocated memory.}
 \label{tab:efficiency}
\end{table}

\paragraph{Readout sensitivity.}
With total width 64 fixed, Table~\ref{tab:sensitivity} varies the number of point-readout heads. Parameter count is unchanged, isolating how the feature width is partitioned. The primary model retains its predefined eight-head configuration; the sensitivity study uses seed 2023 and validation-selected checkpoints.

\begin{table}[t]
 \centering\small\setlength{\tabcolsep}{3pt}
 \begin{tabular}{crrrr}
\toprule
\Lang{头数}{Heads} & Val MSE & Val MAE & Test MSE & Test MAE \\
\midrule
4 & $0.422429$ & $0.427810$ & $0.326837$ & $0.370128$ \\
8 & $0.446899$ & $0.436511$ & $0.336727$ & $0.373539$ \\
16 & $0.425259$ & $0.428136$ & $0.325889$ & $0.369931$ \\
\bottomrule
\end{tabular}

 \caption{Point-readout head sensitivity on ETTh1$(.10,16)$, width 64, seed 2023 and twenty shared test masks.}
 \label{tab:sensitivity}
\end{table}
